\documentclass[11pt]{article}
\usepackage{amsmath,amssymb,amsthm,hyperref}
\newtheorem{theorem}{Theorem}
\newtheorem{lemma}{Lemma}
\newcommand{\E}{\mathbb E}
\newcommand{\F}{\mathbb F}
\newcommand{\OO}{\mathcal A}
\newcommand{\HH}{\mathsf H}
\newcommand{\II}{\mathsf I}
\title{An exact configuration count and polynomial moment\\for finite-alphabet graph intertwiners}
\author{Independent entropy-count audit}
\date{November 15, 2026}
\begin{document}
\maketitle

\paragraph{Scope.}
This note proves the configuration-model counting and entropy step proposed
in \texttt{partitioning-multiclass-coarse-root.txt}. It does not prove
partition rigidity or a recovery theorem. All polynomial exponents below
are explicit but deliberately large. No historical novelty is claimed for
the method of types, symmetrization, or factorial estimates.

Let the planted classes be $V_1,\ldots,V_k$, of sizes $n_i>0$ and total
size $N$. Fix nonnegative integer degrees $d_{ij}$ satisfying
\[
 n_i d_{ij}=n_jd_{ji}\quad(i\ne j),\qquad
 a_i:=d_{ii}\geq2,\qquad n_i a_i\ \hbox{even}.
\]
Write $d_i=\sum_jd_{ij}$ and $M_{ij}=n_i d_{ij}=n_jd_{ji}$ for $i<j$;
write $M_{ii}=n_i a_i$. A vertex of class $i$ has $d_{ij}$ individually
labeled stubs of type $j$. Independently for all factors, pair the internal
stubs uniformly, and match the two sides of every cross factor uniformly.
Loops and multiple edges are retained in this model. An internal loop
contributes two to the diagonal adjacency entry, consistently with its two
stubs and the degree row sum.

Fix a prime $p$, an integer $r\geq1$, the additive alphabet
$\OO=\F_p^r$ of size $q=p^r$, and any matrix $R\in\F_p^{r\times r}$.
A labeling is a row-vector matrix $Y\in\F_p^{N\times r}$. Let
\[
 Z_R(A)=|\{Y:AY=YR\}|.
\]
The entries of $R$ need not be degree counts for this moment theorem.

\begin{theorem}[Polynomial moment with an explicit exponent]
Suppose $k\leq K$, $r\leq K$ and $d_i\leq D$ for every $i$. Put
\[
 q_*:=p^K,\quad S_*:=Kq_*^D,\quad
 F_*:=K+\frac{K(K+3)}2q_*^2,\quad B_*:=S_*+F_*.
 \tag{1}
\]
For every admissible configuration-model size,
\[
 \E Z_R(A)\leq(3D)^{F_*}(N+1)^{B_*}.                 \tag{2}
\]
The bound is independent of $R$ and of the class proportions. If all
factor graphs are conditioned to be simple and the fixed frame is
connected, with fixed positive class proportions, then
\[
 \E[Z_R(A)\mid\text{all factors simple}]
 \leq C(K,D,p)(N+1)^{B_*}                            \tag{3}
\]
along all sufficiently large admissible sizes. The constant can be chosen
uniformly over the finite connected frame family with $k\leq K,d_i\leq D$.
In particular $p=2$ is allowed, so the exponent can depend only on $K,D$.
\end{theorem}

\paragraph{Ordered stars and compatibility.}
Fix a total order on $\OO$. A class-$i$ ordered star is
\[
 s=(x,(z_{j,\ell})_{1\leq j\leq k,\ 1\leq\ell\leq d_{ij}}),
 \qquad x,z_{j,\ell}\in\OO.
\]
Its root label is $x$, and its neighbor labels are ordered by the stub
labels. Retain only allowed stars satisfying
\[
 \sum_{j,\ell}z_{j,\ell}=xR.                         \tag{4}
\]
Because $a_i\geq1$, one internal neighbor label is determined by the root
and all other labels, so the allowed-star set $\mathcal S_i$ has exactly
$q^{d_i}$ elements. Let $T_i(s)$ count the vertices with ordered star $s$.
Its entries are nonnegative integers summing to $n_i$. Define root counts
\[
 t_i(x)=\sum_{s:\operatorname{root}(s)=x}T_i(s)
\]
and aggregated directed-stub counts
\[
 m_{ij}(x,y)=\sum_sT_i(s)
       \sum_{\ell=1}^{d_{ij}}
       \boldsymbol1\{\operatorname{root}(s)=x,z_{j,\ell}=y\}.
 \tag{5}
\]
A realizable type necessarily satisfies
\[
 m_{ii}(x,y)=m_{ii}(y,x),\quad m_{ii}(x,x)\text{ even},\qquad
 m_{ij}(x,y)=m_{ji}(y,x)\quad(i\ne j).                 \tag{6}
\]
Only types satisfying (6) need be included below. Zero-degree factors
are omitted. For an even nonnegative integer $u$, put
\[
 P(u)=(u-1)!!=\frac{u!}{2^{u/2}(u/2)!},\qquad P(0)=1.
\]

\begin{lemma}[Exact factorial formula]
For any compatible collection $T=(T_i)$, its exact contribution to
$\E Z_R(A)$ is
\begin{align}
 W(T)={}&\prod_i\frac{n_i!}{\prod_{s\in\mathcal S_i}T_i(s)!}
 \prod_i\frac{
     \prod_{x<y}m_{ii}(x,y)!\prod_xP(m_{ii}(x,x))}{P(M_{ii})}
 \nonumber\\[-2pt]
 &\hspace{1cm}\times\prod_{i<j:\,M_{ij}>0}
      \frac{\prod_{x,y}m_{ij}(x,y)!}{M_{ij}!}.
 \tag{7}
\end{align}
The expected moment is exactly the sum of (7) over compatible types.
\end{lemma}
\begin{proof}
There are $n_i!/\prod_sT_i(s)!$ assignments of ordered stars to the
labeled vertices of class $i$. This already chooses their root labels;
there is no additional root-label multinomial factor.

For a fixed such assignment, an internal stub with root label $x$ and
target label $y$ must be paired to a stub in the reversed category. For
$x<y$ there are $m_{ii}(x,y)!$ bijections between these categories. For
$x=y$ there are $P(m_{ii}(x,x))$ pairings of that category. The total
number of internal pairings is $P(M_{ii})$. Each cross category $(x,y)$
on side $i$ must match its reversed category on side $j$, giving
$\prod_{x,y}m_{ij}(x,y)!$ compatible bijections out of $M_{ij}!$.
Independence gives (7). Conversely a configuration and a root labeling
determine every ordered star uniquely, including at loops. Thus every
pair contributing to the moment is counted exactly once.
\end{proof}

\paragraph{Normalized laws.}
Let $\rho_i=T_i/n_i$ be the ordered-star law and $\mu_i=t_i/n_i$ its
root law. Put $\nu_{ii}=m_{ii}/M_{ii}$ and
$\nu_{ij}=m_{ij}/M_{ij}$ for active cross factors. Internal pair laws
are symmetric with both marginals $\mu_i$. Each cross pair law has
marginals $\mu_i,\mu_j$ and $\nu_{ji}$ is its reversal. These marginal
claims follow from (5), (6), and the balance identities for $M_{ij}$.
Define, using natural logarithms,
\[
 \Phi(\rho)=\sum_i n_i\HH(\rho_i)
 -\sum_i\frac{n_i a_i}{2}\HH(\nu_{ii})
 -\sum_{i<j:\,M_{ij}>0}M_{ij}\HH(\nu_{ij}).           \tag{8}
\]

\begin{lemma}[Uniform finite-size bound for one type]
With $F_*$ as in (1),
\[
 W(T)\leq\exp\{\Phi(\rho)+F_*[\log(DN+1)+1]\}.      \tag{9}
\]
\end{lemma}
\begin{proof}
Set $f(u)=u\log u-u$ with $f(0)=0$. Integral bounds for the factorial
give, for every integer $u\geq0$,
\[
 0\leq\log(u!)-f(u)\leq\log(u+1)+1.
\]
For even $u$, the displayed formula for $P(u)$ consequently gives
\[
 |\log P(u)-\tfrac12f(u)|\leq2[\log(u+1)+1].
\]
All counts are at most $DN$. The root multinomial contributes at most
$n_i\HH(\rho_i)+[\log(DN+1)+1]$; here the nonnegative errors from
the denominator factorials can simply be dropped for an upper bound.
The main terms of a cross factor are exactly
$-M_{ij}\HH(\nu_{ij})$, with upper error at most
$q^2[\log(DN+1)+1]$.

For an internal factor, symmetry of $m_{ii}$ makes its main term exactly
$-M_{ii}\HH(\nu_{ii})/2$. Its upper error is at most
\[
 \left[\frac{q(q-1)}2+2q+2\right][\log(DN+1)+1]
 \leq2q^2[\log(DN+1)+1],
\]
where $q\geq2$. Summing over at most $k$ root factors, $k$ internal
factors and $k(k-1)/2$ cross factors yields the conservative coefficient
$k+k(k+3)q^2/2\leq F_*$. All estimates include zero counts uniformly.
\end{proof}

\begin{lemma}[Slot symmetrization and weighted entropy cancellation]
Every compatible type satisfies $\Phi(\rho)\leq0$.
\end{lemma}
\begin{proof}
Average each $\rho_i$ over all permutations of its slots within each
neighbor block type, calling the resulting probability law $\bar\rho_i$.
This preserves the root law, all aggregated pair laws, and (4). Concavity
of entropy implies $\HH(\rho_i)\leq\HH(\bar\rho_i)$, hence
$\Phi(\rho)\leq\Phi(\bar\rho)$. The averaged probabilities need not be
integer empirical types; only an entropy inequality is being used.

Work now under these averaged laws. Write $X_i$ for the root, $A_i$ for
any internal slot, and $Z_{ij}$ for any type-$j$ slot. Delete one internal
slot. Equation (4) determines its value from all retained coordinates, so
\[
 \HH(\bar\rho_i)=\HH(\mu_i)+
    \HH(\text{retained neighbors}\mid X_i).
\]
Let $J_i\geq0$ be the conditional total correlation of the retained
neighbors, that is,
\[
 J_i=(a_i-1)\HH(A_i\mid X_i)
       +\sum_{j\ne i}d_{ij}\HH(Z_{ij}\mid X_i)
       -\HH(\text{retained neighbors}\mid X_i).
\]
Nonnegativity follows from conditional entropy subadditivity. Thus
\begin{align*}
 \HH(\bar\rho_i)
 ={}&(2-d_i)\HH(\mu_i)+(a_i-1)\HH(\nu_{ii})
       +\sum_{j\ne i}d_{ij}\HH(\nu_{ij})-J_i.
\end{align*}
Use
\[
 \HH(\nu_{ii})=2\HH(\mu_i)-\II_{ii},\qquad
 \HH(\nu_{ij})=\HH(\mu_i)+\HH(\mu_j)-\II_{ij},
\]
where the $\II$'s are mutual informations. Each cross factor is counted
twice in the star bounds and once in (8). Since $n_i d_{ij}=n_jd_{ji}$,
every root-entropy coefficient cancels exactly, leaving the identity
\[
 \Phi(\bar\rho)=
 -\sum_i n_i(\tfrac{a_i}{2}-1)\II_{ii}
 -\sum_{i<j:\,M_{ij}>0}M_{ij}\II_{ij}
 -\sum_i n_iJ_i\leq0.                               \tag{10}
\]
The last inequality uses $a_i\geq2$. No equality of block sizes, odd
characteristic, or division by a degree was used.
\end{proof}

\begin{proof}[Proof of the moment bound]
There are at most
\[
 \prod_i(n_i+1)^{|\mathcal S_i|}
 \leq(N+1)^{Kq_*^D}=(N+1)^{S_*}
\]
collections of star counts, before any compatibility constraints.
By (9)--(10), each contributes at most
$\exp\{F_*[\log(DN+1)+1]\}$. Since $DN+1\leq D(N+1)$ and $e<3$,
summing proves (2).
\end{proof}

\paragraph{Simplicity conditioning, with the needed classical estimate.}
For completeness, the positive limiting probability of simplicity can be
obtained by the following direct bad-pattern count. In an internal
$a$-regular pairing on $n$ vertices, count loops and unordered pairs of
parallel edges as elementary bad patterns. Their expected counts tend to
\[
 \frac{a-1}{2}\quad\hbox{and}\quad\frac{(a-1)^2}{4},
\]
respectively. In a cross $(b,c)$-biregular matching they tend to no loops
and $(b-1)(c-1)/2$ unordered pairs of parallel edges. To check these
values, choose the involved vertex or vertex pair and the two stubs on
each relevant side, then divide by the total pairing or bijection count.

For every fixed factorial moment order, collections of bad patterns on
disjoint vertices give the product of their limiting means. Consistent
overlapping collections have a bounded graph with more matched edges
than involved vertices and contribute $O(1/\min_i n_i)$; inconsistent
collections contribute zero. This follows by choosing its vertices in
$O(N^v)$ ways and prescribing its $e$ matched pairs with probability
$O(N^{-e})$, with bounded degrees and fixed positive size ratios.
Consequently the sum of bad patterns has factorial moments tending to
those of a Poisson variable with the sum of these means. Applying the
finite Bonferroni inequalities to its zero event, then letting their
order grow, gives simplicity probability tending to the exponential
of the negative sum of means. Simplicity is exactly the absence of
these bad patterns. Distinct factors are independent.

In particular, the limiting exponent for all factors is
\begin{align*}
 \Lambda={}&\sum_i\left[\frac{a_i-1}{2}+\frac{(a_i-1)^2}{4}\right]
  +\sum_{i<j:\,M_{ij}>0}\frac{(d_{ij}-1)(d_{ji}-1)}{2}\\
 \leq{}&\frac{K(D-1)}2+\frac{K^2(D-1)^2}{4}\leq K^2D^2.
\end{align*}
Thus for all sufficiently large admissible sizes the conditioning costs
at most $2e^{K^2D^2}$, so one may use
$C=2\cdot3^{K^2D^2}(3D)^{F_*}$ in (3) eventually. Each simple internal
graph has exactly $(a_i!)^{n_i}$ labeled-stub pairings; each simple cross
graph has exactly $(d_{ij}!)^{n_i}(d_{ji}!)^{n_j}$ labeled-stub
matchings. Hence conditioning produces precisely the intended product
of uniform simple factors, not a uniformity assumption about multigraphs.

For fixed $K,D$ there are finitely many possible degree frames. Connectedness
and balance also give $\min_i n_i\geq N/[KD^{K-1}]$, since a path of at
most $K-1$ frame edges compares any two class sizes by a factor at most
$D^{K-1}$. Therefore the eventual positive simplicity bound can be chosen
uniformly over this finite family and all its admissible size sequences.

\paragraph{Characteristic two and one-hot partitions.}
The vectors $e_1,\ldots,e_r$ remain distinct in $\F_2^r$. For them,
an integer equitable equation reduces to $AY=YR$ over this vector space.
Deleting one internal neighbor in (4) always uses coefficient one, so
neither this step nor any entropy calculation requires an odd prime.
The number of all modular intertwiners has the polynomial bound with
$p=2$, and verified one-hot integer solutions are a subset. The estimate
does not assert that those integer solutions are unique or are only
coarsenings of the planted partition. Those are separate rigidity tasks.

\end{document}
